\SourceLeaf{286}%
% A544-P0286-D01
numeris reliquis (formae $8k+5,8k+7$), si qui inter factores numeri $A$ habentur, prout $S$ est quadratum vel non quadratum. Hinc per ratiocinia prorsus similia ut in art. praec. a serie $f1,f2,f3,f4,f5$ etc., puta $2,2,4,2,0,2$ etc. deinceps ad alias continuo longius \emph{constantes} progrediemur, quarum valores \emph{medii} sint deinceps $\frac\pi{\surd2}$, $\frac\pi{\surd2}\pdot\frac23$, $\frac\pi{\surd2}\pdot\frac23\pdot\frac65$, $\frac\pi{\surd2}\pdot\frac23\pdot\frac65\pdot\frac87$ etc.; progrediemur ita, ut deducamur ad aequationem
\[
2=\frac\pi{\surd2}\pdot\frac23\pdot\frac65\pdot\frac87\pdot\frac{10}{11}\pdot\frac{14}{13}\pdot\frac{16}{17}\pdot\frac{18}{19}\ldots
\]
ubi denominatores constituunt seriem naturalem numerorum primorum, numeratores vero unitate minores sunt, quoties denominatores sunt formae $8k+1$, vel $8k+3$, contra unitate maiores, quoties denominatores sunt formae $8k+5$ vel $8k+7$.
\BookRule{.30\textwidth}

% A544-P0286-D02
\clearpage
\renewcommand{\Article}[1]{\par\medskip\Needspace{4\baselineskip}\pdfbookmark[1]{Art. #1}{nexus-second-prior-art-#1}\centerline{\textbf{#1.}}\nobreak\smallskip}
\pdfbookmark[0]{De nexu inter multitudinem classium. II}{nexus-second-prior}
\begin{center}
[II.]\par\medskip
\Large\bfseries DE NEXU INTER MULTITUDINEM CLASSIUM, IN QUAS FORMAE BINARIAE SECUNDI GRADUS DISTRIBUUNTUR, EARUMQUE DETERMINANTEM.\par\medskip
\normalsize COMMENTATIO PRIOR\par\medskip
SOCIETATI REGIAE EXHIBITA 1837\ .\ .\ .
\end{center}
\BookRule{.18\textwidth}

% A544-P0286-D03
\Article{1}
Triginta sex elapsi sunt anni, ex quo principia nexus mirabilis in hac commentatione tractandi detecta sunt, uti iam in fine \emph{Disquisitionum arithmeticarum} annuntiatum est. Sed aliae occupationes per longum tempus ab hac scrutatione detraxerant, donec recentiori tempore ad eam reverti et per novas curas eam ampliare contigerit. Attamen quum ambitus huius novae Arithmeticae Sublimioris partis limites unius commentationis transgrediatur, haecce prior formis determinantium negativorum dicata erit: formae autem determinantium positivorum, quae tractationem prorsus peculiarem requirunt, commentationi alteri reservatae manebunt.

\SourceLeaf{287}%
% A544-P0287-D01
\Article{2}
Ad propositum nostrum opus erit theoremate per se quidem arithmetico, cuius tamen indolem commodius et clarius per considerationes in forma geometrica exhibendas ob oculos ponere licet.

% A544-P0287-D02
Proposita in plano indefinito figura per lineam qualemcunque terminata, illius area approximative assignari poterit, si plano in quadrata dispertito multitudo tum eorum quae integra sunt intra figuram, tum eorum quae ambitus figurae secat, numeretur, manifestoque area iusto minor vel maior prodibit, prout quadrata posteriora vel omittuntur vel prioribus adnumerantur: si vero quadrata posteriora in limine sita, ad normam qualiscunque principii, partim excludere partim adnumerare placuerit, error modo positivus modo negativus esse poterit, necessario tamen minor quam aggregatum cunctorum quadratorum in limine. Quo minora quadrata accipiantur, eo exactius hoc modo area determinabitur, talemque approximationem in infinitum producere sive quadrata tam parva accipere licebit, ut error quavis quantitate data minor evadat. Quod quamquam iam per se evidens esse videatur, tamen demonstratione rigorosa munire non aspernabimur.

% A544-P0287-D03
Bina quadrata vel unum punctum angulare, vel duo, vel nullum commune habere possunt; in casu primo et secundo contigua, in tertio disiuncta dicentur. Manifesto quadrata, quae omnia inter se contigua sint, quaterna tantum exstant, adeoque inter quina quadrata diversa duo ad minimum disiuncta inveniri debent. Iam quum distantia inter duo puncta in quadratis disiunctis sita nequeat esse minor quam latus quadratorum, quod per $a$ designabimus, patet, si punctum a quocunque alicuius quadrati loco profectum deinceps quadratum secundum, tertium, quartum traiecerit, tandem ad quintum pervenerit, longitudinem viae certe non esse minorem quam $a$. Et quum simili ratione si linea continuo alia quadrata permeat, pars inter quadratum quintum et nonum, nec non inter nonum et decimum tertium etc. non possit esse minor quam $a$, facile colligimus, lineam quamcunque in se ipsam redeuntem, quae omnino $n$ quadrata diversa attigerit, certo non posse esse minorem quam $\frac{(n-4)a}4$. Vice versa itaque linea clausa, cuius longitudo est $=l$, certo plura quam $4+\frac{4l}a$ quadrata diversa attigisse non potest. Quorum area $=4aa+4al$ quum decrescente $a$ in infinitum quavis quantitate data minor fieri possit, idem a potiori valebit de errore quadraturae de qua supra diximus.

\SourceLeaf{288}%
% A544-P0288-D01
\Article{3}
Principium admissionis vel exclusionis quadratorum in limite figurae positorum multis modis diversis condi posset: simplicissimum tamen videtur, tantummodo situm centri cuiusque quadrati respicere, ita ut admittantur quadrata, quorum centra sunt intra figuram, excludantur ea, quorum centra sunt extra figuram, denique arbitrio relinquatur, utrum centra, quae forte in peripheria ipsa sunt, interioribus vel exterioribus adnumerare malimus. Loco centrorum etiam quaevis alia puncta in singulis quadratis similiter sita adoptare possemus.

% A544-P0288-D02
Hoc pacto res eo redit, ut in plano puncta aequidistantia et in rectis aequidistantibus ita disseminata concipiamus, ut quadrata offerant: quo facto per theorema art. praec. affirmare possumus, multitudinem punctorum in figura contentorum in quadratum distantiae binorum punctorum proximorum multiplicatam areae figurae quam prope velis aequalem evadere, si modo distantia ista satis parva accipiatur, sive ad instar vulgaris loquendi modi, productum illud aream exhibere, si distantia sit infinite parva.

% A544-P0288-D03
\Article{4}
Curva per aequationem inter coordinatas orthogonales $p,q$ hancce
\[
app+2bpq+cqq=1
\]
expressa, est sectio conica, et quidem ellipsis, si $a,c$ atque $ac-bb$ sunt quantitates positivae: area hac ellipsi circumscripta invenitur $\pdot=\frac\pi{\surd(ac-bb)}$. Valor quantitatis $app+2bpq+cqq$ extra ellipsem ubique fit maior quam $1$, intra ellipsem minor quam $1$, negativus nullibi.

% A544-P0288-D04
Concipiatur .systema punctorum per planum, in quo ellipsis sita est, ita disseminatorum, ut forment quadrata, quorum latera $=\lambda$ axibus coordinatarum sint parallela, ubi nihil refert, utrum initium coordinatarum sive centrum ellipsis cum aliquo horum punctorum coincidat necne. Sit multitudo punctorum intra ellipsem, adnumeratis si quae sunt in ipsa peripheria, $=m$, eritque per theorema art. praec. $\frac\pi{\surd(ac-bb)}$ limes quantitatis $m\lambda\lambda$, ad quem quam prope velis accedit, decrescente $\lambda$ in infinitum.

% A544-P0288-D05
Si initium coordinatarum cum aliquo systematis puncto coincidere supponimus, statuendo $p=\lambda x,q=\lambda y$, manifesto pro singulis punctis systematis $x$ et $y$ erunt numeri integri, et vice versa quaevis combinatio valorum integrorum
\SourceLeaf{289}%
% A544-P0289-D01
quantitatum $x,y$ respondebit alicui systematis puncto. Hinc numerus $m$ nihil aliud est, nisi multitudo omnium combinationum valorum integrorum quantitatum $x,y$, pro quibus $F$ non fit maior quam $M$, si brevitatis caussa functionem, seu formam secundi ordinis $axx+2bxy+cyy$ per $F$, atque quantitatem $\frac1{\lambda\lambda}$ per $M$ denotamus. Determinans huius formae est $bb-ac$, pro quo scribemus $-D$. Hoc pacto theorema nostrum iam ita enunciandum erit.

% A544-P0289-D02
\textsc{Theorema I.}\quad\emph{Multitudo $m$ omnium combinationum valorum integrorum indeterminatarum $x,y$, pro quibus valor formae determinantis negativi $-D$ limitem $M$ non egreditur, fit $=\frac{\pi M}{\surd D}$, proxime quidem, sed approximatione in infinitum crescente, dum $M$ crescit in infinitum. Vix erit monendum, approximationem infinitam hic (et perinde in sequentibus) non ita intelligendam, ac si differentia inter $\frac{\pi M}{\surd D}$ et $m$ ipsa in infinitum decrescat, sed ratio inter has quantitates ad aequalitatem in infinitum appropinquabit, sive $\frac{\pi M}{m\surd D}-1$ in infinitum decrescet.}

% A544-P0289-D03
\Article{5}
Ad dinumerationem reapse efficiendam ita procedi potest, ut pro singulis valoribus integris ipsius $x$ inter limites $-\surd\frac{cM}D$ atque $+\surd\frac{cM}D$ sitis bini valores ipsius $y$ aequationi $F=M$ respondentes computentur, unde multitudo integrorum inter hos iacentium sponte habetur. Quum haec multitudo eadem sit pro valoribus oppositis ipsius $x$, laboris dimidia fere parte liberamur. Res ita quoque perfici potest, ut valores ipsius $x$ dinumerentur singulis valoribus ipsius $y$ inter limites $-\surd\frac{aM}D$ atque $+\surd\frac{aM}D$ respondentes. Per combinationem idoneam utriusque methodi labor amplius sublevari potest, quod tamen fusius hic non exsequimur: sufficiat de casu simplicissimo quaedam adiungere.

% A544-P0289-D04
Sit forma $F=xx+yy$, sive curva circulus, designentque $r,r',r'',r'''\ldots r^{(r)}$ numeros integros proxime minores quam
\[
\surd M,\quad\surd(M-1),\quad\surd(M-4),\quad\surd(M-9)\ldots\surd(M-rr)
\]
vel si quae inter has quantitates sunt integri, hos ipsos. Tunc erit multitudo quaesita
\[
\begin{aligned}
m&=2r+1+2(2r'+1)+2(2r''+1)+2(2r'''+1)+\text{etc.}+2(2r^{(r)}+1)\\
&=1+4r+4r'+4r''+4r'''+\text{etc.}+4r^{(r)}
\end{aligned}
\]

% A544-P0289-D05
Expeditius autem idem assequimur, denotando per $q$ integrum proxime
\SourceLeaf{290}%
% A544-P0290-D01
minorem quam $\surd\tfrac12 M$ (vel hanc quantitatem ipsam, si fit numerus integer) adiumento formulae
\[
m=4qq+1+4r+8(r^{(q+1)}+r^{(q+2)}+r^{(q+3)}+\text{etc.}+r^{(r)})
\]
Hoc modo eruta sunt sequentia.
\[
\begin{array}{r|r||r|r||r|r}
M&m&M&m&M&m\\\hline
100&317&1000&3149&10000&31417\\
200&633&2000&6293&20000&62845\\
300&949&3000&9425&30000&94237\\
400&1257&4000&12581&40000&125629\\
500&1581&5000&15705&50000&157093\\
600&1885&6000&18853&60000&188453\\
700&2209&7000&21993&70000&219901\\
800&2521&8000&25137&80000&251305\\
900&2821&9000&28269&90000&282697\\
1000&3149&10000&31417&100000&314197
\end{array}
\]

% A544-P0290-D02
\Article{6}
Theoremati art. 4 maiorem generalitatem conciliamus sequenti modo.

% A544-P0290-D03
\textsc{Theorema II.}\quad\emph{Si non omnes combinationes valorum integrorum quantitatum $x,y$ pro quibus $F$ non egreditur valorem $M$, colligendae sunt, sed tantummodo per saltus, puta eae, ubi $x$ congruus est numero dato $G$ secundum modulum datum $g$, atque $y$ congruus numero dato $H$ secundum modulum datum $h$, harum combinationum multitudo $m'$ exprimetur proxime per $\frac{\pi M}{gh\surd D}$, approximatione in infinitum aucta, dum $M$ in infinitum crescet.}

% A544-P0290-D04
Revera statuendo $x=gx'+G,y=hy'+H$, patet, $m'$ esse multitudinem omnium combinationum valorum integrorum quantitatum $x',y'$, pro quibus
\[
agg(x'+\tfrac Gg)^2+2bgh(x'+\tfrac Gg)(y'+\tfrac Hh)+chh(y'+\tfrac Hh)^2
\]
valorem $M$ non egrediatur. Manifesto igitur si in plano systema punctorum perinde quidem ut in art. 4 disseminatum supponimus, attamen ita ut non initium coordinatarum sed punctum, cuius coordinatae sunt $p=\frac{G\lambda}g,q=\frac{H\lambda}h$, cum aliquo systematis puncto coincidat, $m'$ exprimet multitudinem punctorum intra 
